\subsection{泛系数公式}

在本节中，我们考虑两个 A-模复形 (C, d) 和 (C', d')。考虑如下典范正合序列：

\[
0 \rightarrow \mathrm{Z} (\mathrm{C}) \stackrel {{j}} {{\rightarrow}} \mathrm{C} \stackrel {{\delta}} {{\rightarrow}} \mathrm{B} (\mathrm{C}) (- 1) \rightarrow 0,\tag{I}
\]

\((\mathrm{II}_{p})\)

\[
0 \to \mathrm{B} _ {p} (\mathrm{C}) \stackrel {{i}} {{\to}} \mathrm{Z} _ {p} (\mathrm{C}) \stackrel {{p}} {{\to}} \mathrm{H} _ {p} (\mathrm{C}) \to 0;
\]

我们从 δ 推出一个 k-同态

\[
\mathrm{H} (\text { Homgr } (\delta , 1)): \mathrm{H} (\text { Homgr } _ {\mathbf {A}} (\mathrm{B} (\mathbf {C}), \mathbf {C} ^ {\prime})) (1) \rightarrow \mathrm{H} (\text { Homgr } _ {\mathbf {A}} (\mathbf {C}, \mathbf {C} ^ {\prime}));
\]

由 \((\mathrm{II}_{p})\) 推出连接同态：

\[
\delta \left(\mathrm{II} _ {p}, \mathrm{H} ^ {q} \left(\mathrm{C} ^ {\prime}\right)\right): \operatorname{Hom} _ {\mathrm{A}} \left(\mathrm{B} _ {p} (\mathrm{C}), \mathrm{H} ^ {q} \left(\mathrm{C} ^ {\prime}\right)\right)\rightarrow \operatorname{Ext} _ {\mathrm{A}} ^ {1} \left(\mathrm{H} _ {p} (\mathrm{C}), \mathrm{H} ^ {q} \left(\mathrm{C} ^ {\prime}\right)\right).
\]

由此，通过过渡到乘积，得到 k-模的同态：

\[
\varphi^ {n}: \operatorname{Homgr} _ {\mathbf {A}} ^ {n} (\mathrm{B} (\mathrm{C}), \mathrm{H} (\mathrm{C} ^ {\prime})) \rightarrow \prod_ {p + q = n} \operatorname{Ext} _ {\mathbf {A}} ^ {1} \left(\mathrm{H} _ {p} (\mathrm{C}), \mathrm{H} ^ {q} (\mathrm{C} ^ {\prime})\right)
\]

还有典范同态（X，第 82 页）

\[
\lambda^ {n} (\mathrm{B} (\mathrm{C}), \mathrm{C} ^ {\prime}): \mathrm{H} ^ {n} \left(\operatorname{Homgr} _ {\mathrm{A}} (\mathrm{B} (\mathrm{C}), \mathrm{C} ^ {\prime})\right)\rightarrow \operatorname{Homgr} _ {\mathrm{A}} ^ {n} (\mathrm{B} (\mathrm{C}), \mathrm{H} (\mathrm{C} ^ {\prime})) .
\]

在这些记号下：

定理3. --- 假设 A-模 B(C) 和 Z(C) 是投射的。那么对每个 n，存在唯一的 k-模同态

\[
\beta^ {n}: \prod_ {p + q = n - 1} \operatorname{Ext} _ {\mathbf {A}} ^ {1} \left(\mathrm{H} _ {p} (\mathbf {C}), \mathrm{H} ^ {q} \left(\mathbf {C} ^ {\prime}\right)\right)\rightarrow \mathrm{H} ^ {n} \left(\operatorname{Homgr} _ {\mathbf {A}} \left(\mathrm{C}, \mathrm{C} ^ {\prime}\right)\right)
\]

从而使得图表交换

\[
\begin{array}{c} \operatorname{Homgr} _ {\mathbf {A}} ^ {n - 1} (\mathrm{B} (\mathrm{C}), \mathrm{H} (\mathrm{C} ^ {\prime})) \xleftarrow {\lambda^ {n - 1} (\mathrm{B} (\mathrm{C}) , \mathrm{C} ^ {\prime})} \mathrm{H} ^ {n - 1} (\operatorname{Homgr} _ {\mathbf {A}} (\mathrm{B} (\mathrm{C}), \mathrm{C} ^ {\prime})) \\ \prod_ {p + q = n - 1} \operatorname{Ext} _ {\mathbf {A}} ^ {1} (\mathrm{H} _ {p} (\mathrm{C}), \mathrm{H} ^ {q} (\dot {\mathrm{C}} ^ {\prime})) \xrightarrow {\beta^ {n}} \mathrm{H} ^ {n} (\operatorname{Homgr} _ {\mathbf {A}} (\mathrm{C}, \mathrm{C} ^ {\prime})). \end{array}
\]

分次 k-模的序列

\[
\begin{array}{r l} 0 \longrightarrow & \prod_ {p + q = n - 1} \operatorname{Ext} _ {\mathrm{A}} ^ {1} \left(\mathrm{H} _ {p} (\mathrm{C}), \mathrm{H} ^ {q} (\mathrm{C} ^ {\prime})\right) \xrightarrow {\beta^ {n}} \mathrm{H} ^ {n} \big (\mathrm{Homg} _ {\mathrm{A}} (\mathrm{C}, \mathrm{C} ^ {\prime}) \big) \\ & \xrightarrow {\lambda^ {n} (\mathrm{C} , \mathrm{C} ^ {\prime})} \prod_ {p + q = n} \mathrm{Hom} _ {\mathrm{A}} \left(\mathrm{H} _ {p} (\mathrm{C}), \mathrm{H} ^ {q} (\mathrm{C} ^ {\prime})\right) \to 0 \end{array}\tag{11}
\]

是正合的。

注记. --- 可以证明一个类似的命题，假设 \(\mathbf{B}_{n}(\mathbf{C}')\) 和 \(\mathbf{C}'/\mathbf{B}_{n}(\mathbf{C}')\) 对每个 n 都是内射的。

为简单起见，设 \(\mathbf{B} = \mathbf{B}(\mathbf{C}), \mathbf{Z} = \mathbf{Z}(\mathbf{C}), \mathbf{H} = \mathbf{H}(\mathbf{C})\) 和 \(\mathbf{H}' = \mathbf{H}(\mathbf{C}')\) 。由于 B 是投射的，由 (I) 可导出一个正合序列

\[
\begin{array}{r l}(1 2) \quad&0 \rightarrow \operatorname{Homgr} _ {\mathbf {A}} (\mathrm{B}, \mathrm{C} ^ {\prime}) (1) \xrightarrow {\operatorname{Homgr} (\delta , 1)} \operatorname{Homgr} _ {\mathbf {A}} (\mathrm{C}, \mathrm{C} ^ {\prime})\\&\xrightarrow {\operatorname{Homgr} (j , 1)} \operatorname{Homgr} _ {\mathbf {A}} (\mathrm{Z}, \mathrm{C} ^ {\prime}) \rightarrow 0.\end{array}
\]

引理2. --- 连接同态

\[
\mathrm{H} ^ {n} \left(\operatorname{Homgr} _ {\mathbf {A}} \left(\mathrm{Z}, \mathrm{C} ^ {\prime}\right)\right)\rightarrow \mathrm{H} ^ {n} \left(\operatorname{Homgr} _ {\mathbf {A}} \left(\mathrm{B}, \mathrm{C} ^ {\prime}\right)\right)
\]

与 (12) 相关联的连接同态等于 \((-1)^{n+1}\) H(Homgr(i, 1)).

事实上，设 \(a \in \mathbf{Z}^n(\text{Homgr}_A(Z, C'))\) ；这是一个次数递增 \(n\) 的、从 \(Z\) 到 \(C'\) 的复形态射，其值因此位于 \(Z(C')\) 中。由于正合序列 (I) 是分裂的（因为 B 是投射的）， \(a\) 延拓为一个元素 \(b\) ，它属于 \(\text{Homgr}_A^n(C, Z(C'))\) 。根据定义， \(a\) 的类在该连接同态下的像是 \(H^n(\text{Homgr}_A(B, C'))\) 中同态 \(u\) （从 \(B(C)\) 到 \(C'\) ）的类，使得对于 \(x \in C\) ，有

\[
u (d x) = \mathrm{D} b (x) = d ^ {\prime} b (x) - (- 1) ^ {n} b (d x) = (- 1) ^ {n + 1} b (d x) = (- 1) ^ {n + 1} a (d x),
\]

由此得出断言。

因此，与（12）相关联的同调正合列给出如下正合序列

\[
\rightarrow \mathrm{H} ^ {n} (\operatorname{Homgr} _ {\mathrm{A}} (Z, C ^ {\prime})) \xrightarrow {\mathrm{H} (\operatorname{Homgr} (i , 1))} \mathrm{H} ^ {n} (\operatorname{Homgr} _ {\mathrm{A}} (B, C ^ {\prime}))
\]

\[
\xrightarrow {\mathrm{H} (\text { Homgr } (\delta , 1))} \mathrm{H} ^ {n + 1} \left(\text { Homgr } _ {\mathrm{A}} \left(\mathrm{C}, \mathrm{C} ^ {\prime}\right)\right) \xrightarrow {\mathrm{H} (\text { Homgr } (j , 1))} \mathrm{H} ^ {n + 1} \left(\text { Homgr } _ {\mathrm{A}} \left(\mathrm{Z}, \mathrm{C} ^ {\prime}\right)\right)\rightarrow \dots .
\]

此外，由于 Z 是投射的，由 \((II_{p})\) 可导出正合序列

\[
0 \to \mathrm{Hom} _ {\mathbf {A}} (\mathrm{H} _ {p}, \mathrm{H} ^ {\prime q}) \to \mathrm{Hom} _ {\mathbf {A}} (Z _ {p}, \mathrm{H} ^ {\prime q}) \to \mathrm{Hom} _ {\mathbf {A}} (\mathrm{B} _ {p}, \mathrm{H} ^ {\prime q}) \to \mathrm{Ext} _ {\mathbf {A}} ^ {1} (\mathrm{H} _ {p}, \mathrm{H} ^ {\prime q}) \to 0,
\]

由此，通过取乘积，得到正合序列

\[
\begin{array}{r l} 0 \to \operatorname{Homgr} _ {\mathbf {A}} ^ {n} (\mathbf {H}, \mathbf {H} ^ {\prime}) & \to \operatorname{Homgr} _ {\mathbf {A}} ^ {n} (Z, \mathbf {H} ^ {\prime}) \to \operatorname{Homgr} _ {\mathbf {A}} ^ {n} (\mathbf {B}, \mathbf {H} ^ {\prime}) \\ & \to \prod_ {p + q = n} \operatorname{Ext} _ {\mathbf {A}} ^ {1} (\mathbf {H} _ {p}, \mathbf {H} ^ {\prime q}) \to 0. \end{array}
\]

最后，有第 1 号中的典范同态：

\[
\begin{array}{l} \lambda_ {\mathrm{B}} = \lambda (\mathrm{B}, \mathrm{C} ^ {\prime}): \mathrm{H} (\mathrm{Homgr} _ {\mathrm{A}} (\mathrm{B}, \mathrm{C} ^ {\prime})) \to \mathrm{Homgr} _ {\mathrm{A}} (\mathrm{B}, \mathrm{H} ^ {\prime}), \\ \lambda_ {\mathrm{Z}} = \lambda (\mathrm{Z}, \mathrm{C} ^ {\prime}): \mathrm{H} (\mathrm{Homgr} _ {\mathrm{A}} (\mathrm{Z}, \mathrm{C} ^ {\prime})) \to \mathrm{Homgr} _ {\mathrm{A}} (\mathrm{Z}, \mathrm{H} ^ {\prime}), \\ \lambda_ {\mathrm{C}} = \lambda (\mathrm{C}, \mathrm{C} ^ {\prime}): \mathrm{H} (\mathrm{Homgr} _ {\mathrm{A}} (\mathrm{C}, \mathrm{C} ^ {\prime})) \to \mathrm{Homgr} _ {\mathrm{A}} (\mathrm{H}, \mathrm{H} ^ {\prime}), \end{array}
\]

由此得到 X 第 97 页上的行正合图表。

该图表由 \(\lambda\) 同态的构造是交换的。此外，由于复形 B 和 Z 是分裂且投射的， \(\lambda_{\mathrm{B}}\) 和 \(\lambda_{\mathrm{Z}}\) 是双射的（X，第 84 页，推论 1）。由此一方面可推出， \(\lambda_{\mathrm{C}}^{n}\) 是满射的，其核等于 Im H \(^{n}\) (Homgr ( \(\delta\) , 1))，另一方面可推出， \(\varphi^{n-1} \circ \lambda_{\mathrm{B}}^{n-1}\) 是满射的，其核等于 Ker H \(^{n}\) (Homgr ( \(\delta\) , 1))。定理由此立即得出。

推论 1. --- 假设 B(C) 和 Z(C) 是投射的，且 B \(^{n}\) (C') 对每个 n 都是内射的。则正合序列 (11) 是分裂的。

这由定理及以下引理得出：

引理 3. --- 若 B(C) 是投射的且 \(\mathbf{B}_{n}(\mathbf{C}^{\prime})\) 对每个 \(n\) ，典范同态 \(\lambda(\mathbf{C}, \mathbf{C}^{\prime}): \mathrm{H}(\mathrm{Homgr}_{\mathbf{A}}(\mathbf{C}, \mathbf{C}^{\prime})) \to \mathrm{Homgr}_{\mathbf{A}}(\mathrm{H}(\mathbf{C}), \mathrm{H}(\mathbf{C}^{\prime}))\) 有一个 \(k\) -线性截面。

事实上，根据 X，第 35 页，注记 a) 和 b)，存在拟同构

\[
\varphi : \mathrm{C} \rightarrow \mathrm{H} (\mathrm{C}) \quad \text { and } \quad \varphi^ {\prime}: \mathrm{H} (\mathrm{C} ^ {\prime}) \rightarrow \mathrm{C} ^ {\prime}
\]

使得 \(\mathbf{H}(\varphi) = 1_{\mathbf{H}(\mathbf{C})}\) 和 \(\bar{\mathbf{H}} (\varphi^{\prime}) = 1_{\mathbf{H}(\mathbf{C}^{\prime})}\) .

在交换图中

\[
\begin{array}{l} \mathrm{H} (\operatorname{Homgr} _ {\mathbf {A}} (\mathrm{H} (\mathrm{C}), \mathrm{H} (\mathrm{C} ^ {\prime})) \xrightarrow {\mathrm{H} (\operatorname{Homgr} (\varphi , \varphi^ {\prime}))} \mathrm{H} (\operatorname{Homgr} _ {\mathbf {A}} (\mathrm{C}, \mathrm{C} ^ {\prime})) \\ \lambda (\mathrm{H} (\mathrm{C}), \mathrm{H} (\mathrm{C} ^ {\prime})) \Big | _ {\downarrow} \\ \operatorname{Homgr} _ {\mathbf {A}} (\mathrm{H} (\mathrm{C}), \mathrm{H} (\mathrm{C} ^ {\prime})) \xrightarrow {\operatorname{Homgr} (\mathrm{H} (\varphi) , \mathrm{H} (\varphi^ {\prime}))} \operatorname{Homgr} _ {\mathbf {A}} (\mathrm{H} (\mathrm{C}), \mathrm{H} (\mathrm{C} ^ {\prime}))  , \end{array}
\]

\(\lambda(H(C), H(C'))\) 是双射且 Homgr \((H(\varphi), H(\varphi'))\) 是恒等映射，由此得出断言。

Homgr \(^{n-1}\) (Z, H') → Homgr \(^{n-1}\) (B, H') → Φ \(^{n-1}\) → Πp+q=n-1 Ext \(^{1}\) (Hp, H'q) → 0

H \(^{n-1}\) (Homgr(A, C')) → H \(^{n-1}\) (Homgr(A, B, C')) → H \(^{n}\) (Homgr(A, C)) → H \(^{n}\) (Homgr(A, C)) → H \(^{n}\) (Homgr(A, B, C))

0 → Homgr \(^{n}\) (H, H') → Homgr \(^{n}\) (Z, H') → Homgr \(^{n}\) (B, H')

推论 2. --- 若 B(C) 和 Z(C) 是投射的，则对每个 A-模 N 和每个整数 n，存在一个分裂的正合序列：

\[
(1 3) \quad 0 \longrightarrow \operatorname{Ext} _ {\mathbf {A}} ^ {1} \left(\mathrm{H} _ {n - 1} (\mathrm{C}), \mathrm{N}\right) \xrightarrow {\beta^ {n}} \mathrm{H} ^ {n} \left(\operatorname{Homgr} _ {\mathbf {A}} (\mathrm{C}, \mathrm{N})\right) \xrightarrow {\lambda^ {n}} \operatorname{Hom} _ {\mathbf {A}} \left(\mathrm{H} _ {n} (\mathrm{C}), \mathrm{N}\right) \longrightarrow 0.
\]

推论 3（“泛系数公式”）. --- 设 A 是主理想环且 C 是自由的。则对每个 A-模 N 和每个整数 n，存在一个分裂的正合序列 (13)。

事实上，B(C) 和 Z(C) 作为自由模 C 的子模是自由的（VII，§ 3，定理 1 的推论 2）。

推论 4. --- 若 C 右有界且 C 和 H(C) 是投射的，则

\[
\lambda (C, C ^ {\prime}): H \left(\operatorname{Homgr} _ {A} (C, C ^ {\prime})\right)\rightarrow \operatorname{Homgr} _ {A} \left(H (C), H \left(C ^ {\prime}\right)\right)
\]

是双射。

由定理可知，只需证明 B(C) 和 Z(C) 是投射的。现在我们有正合序列

\[
\begin{array}{l} 0 \to \mathrm{B} _ {n} (\mathrm{C}) \to Z _ {n} (\mathrm{C}) \to \mathrm{H} _ {n} (\mathrm{C}) \to 0 \\ 0 \to Z _ {n} (\mathrm{C}) \to \mathrm{C} _ {n} \to \mathrm{B} _ {n - 1} (\mathrm{C}) \to 0, \end{array}
\]

因此 \((\mathbf{B}_{n - 1}(\mathbf{C})\) 是投射的) \(\Rightarrow (\mathbf{Z}_n(\mathbf{C})\) 是投射的) \(\Rightarrow (\mathbf{B}_n(\mathbf{C})\) 是投射的)。我们通过指出 \(\mathbf{B}_n(\mathbf{C}) = 0\) （对足够小的 \(n\) ）来结束。

